\chapter{Knowing that you are right}

\begin{orient}
\textbf{What this chapter is for.} A solution you cannot check is a guess with working attached. This chapter treats verification as a technique with its own cues, its own methods, and its own traps, exactly like substitution or parts. It presents eight checks ordered by cost, from a five-second scaling argument to a full recomputation by an independent route, and it names the specific errors each one is built to catch. At the end you should have a reflex: on finishing any problem, you reach for the cheapest check that could possibly falsify the answer you just wrote, and you run it before you believe yourself.
\end{orient}

\section{Verification is a technique, not a virtue}

Checking your work is usually taught as a moral instruction, which is why it is usually ignored. It is better understood as a decision problem. You have an answer produced by a chain of ten steps, each of which you believe with probability $0.98$; the answer is therefore about $80$ percent likely to be right, which is not good enough. A check that would fail on most wrong answers and passes on this one moves that number a long way for very little effort. The question is never whether to check, but which check has the best ratio of discriminating power to cost.

That framing has two consequences. The first is that a check must be capable of failing. Recomputing the same algebra a second time is not verification, because the error that survived the first pass survives the second; the mind that made the slip is the mind rechecking it. A real check uses different information: a bound, a special case, a numerical value, a symmetry, a second technique. The second consequence is that checks should be run in order of cost, because a cheap check that kills the answer saves you from paying for the expensive one.

\begin{keynote}[The discriminating-check rule]
A check is worth running in proportion to the fraction of plausible wrong answers it would reject. Confirming that a positive integral came out positive is nearly free but rejects only the sign errors; confirming that it agrees with numerical quadrature to six digits rejects essentially everything. Spend the first ten seconds on the checks that cost nothing, and the next two minutes on the one check that would be decisive.
\end{keynote}

\section{The verification ladder}

The eight checks below are arranged by cost. The first four cost seconds and should be automatic; the last four cost minutes and are what you run on an answer that matters. On any given problem two or three of them will apply, and applying two independent ones is worth much more than applying one twice.

\tech{Dimensional and scaling sanity}{core}
\cue An answer containing a parameter that carries a scale: a length $a$, a rate $k$, an upper limit $T$. Also any integral over a scale-invariant domain such as $(0,\infty)$.
\method Substitute $x=\lambda u$ and see how both sides must transform. If the integrand is homogeneous of degree $d$ in $x$ and $a$ jointly and the domain is scale-invariant, then the answer is $a^{\,d+1}$ times a pure number, and the power of $a$ is forced before any computation. Any claimed closed form with the wrong power is wrong.

\begin{worked}
Claim: $\displaystyle\int_{0}^{\infty}\frac{\dd x}{(x^{2}+a^{2})^{2}}=\frac{\pi}{4a^{3}}$ for $a>0$. Substituting $x=au$ gives $\frac{1}{a^{3}}\int_{0}^{\infty}\frac{\dd u}{(u^{2}+1)^{2}}$, so the answer is forced to be $c\,a^{-3}$; the constant $c=\pi/4$ then follows from the single case $a=1$. A claimed answer of $\pi/(4a^{2})$ or $\pi/(2a^{3})$ is refuted by the first half of that argument alone.
\end{worked}

\begin{trap}
Using the scaling argument on a domain that is not scale-invariant. On $[0,1]$ the substitution $x=au$ moves the upper limit, so it yields a relation between two different integrals, not a constraint on one.
\end{trap}

\begin{seealsob}Degenerate and limiting cases; symmetry of the answer.\end{seealsob}

Scaling constrains the shape of the answer. The next check constrains its value, by pushing a parameter to a place where you already know what the answer must be.

\tech{Degenerate and limiting cases}{core}
\cue Any closed form containing a parameter, an index, a dimension, or a count. Also any general formula you have just derived and are about to trust.
\method Set the parameter to $0$, to $1$, or let it tend to $\infty$, and check that the formula degenerates to something you know independently. Prefer the extreme case to the convenient one: at $a=1$ many wrong answers coincide with the right one, whereas at $a\to 0$ and $a\to\infty$ they separate.

\begin{worked}
Claim: $\displaystyle\int_{0}^{1}x^{m}(1-x)^{n}\dd x=\frac{m!\,n!}{(m+n+1)!}$ for non-negative integers $m,n$. At $n=0$ the integral is $\frac{1}{m+1}$ and the formula gives $\frac{m!\,0!}{(m+1)!}=\frac{1}{m+1}$. At $m=n=0$ both sides are $1$. At $m=2$, $n=5$ the formula gives $\frac{2\cdot 120}{40320}=\frac{1}{168}=0.0059524$, and Simpson's rule on the integral returns $0.0059524$.
\end{worked}

\begin{trap}
Choosing a degenerate case that the wrong answer also satisfies. If a suspected error is a missing factor of $\frac12$, then $m=n$ will not expose it; pick the case where the suspected error is maximal.
\end{trap}

\begin{seealsob}Small-case brute force; dimensional and scaling sanity.\end{seealsob}

If a parameter is not available, evaluate. Numerical evaluation is the most powerful cheap check in the book, and the only one that essentially never passes a wrong answer.

\tech{Numerical evaluation to several digits}{core}
\cue Any definite integral, any limit, any series with a claimed closed form, any probability.
\method Get four to six significant figures. Simpson's rule for smooth integrands; the midpoint rule for periodic ones and for anything with an endpoint singularity, since it never evaluates the endpoint; a substitution that removes the singularity when the midpoint rule converges too slowly. For a limit, evaluate at $n=10,10^{2},10^{3},10^{4}$ and read the trend, not the last value.

\begin{worked}
Claim: $\displaystyle\int_{0}^{\pi/2}\frac{\dd x}{1+\sin^{2}x}=\frac{\pi}{2\sqrt2}=1.1107207$. The midpoint rule with only eight subintervals returns $1.1107207$, agreeing to seven decimals, because the integrand extends to a smooth periodic function and the midpoint rule then converges geometrically. Eight function evaluations settle the question.
\end{worked}

\begin{trap}
Accepting agreement to two digits. A wrong answer differing from the right one by a factor of $\frac{\pi^{2}}{10}$ or by $\frac{101}{100}$ passes two digits comfortably. Insist on four, and on more than that when the suspected error is a small rational factor.
\end{trap}

\begin{seealsob}Sign and order of magnitude; small-case brute force.\end{seealsob}

For an antiderivative or a differential equation, there is a check that is not merely strong but complete: run the problem backwards. Differentiation is algorithmic in a way that integration is not, which is exactly what makes it a reliable inverse.

\tech{Differentiate the answer back}{core}
\cue An indefinite integral, a solution to an ODE, an inverse function, an implicit relation, or any claim of the form ``$F$ is an antiderivative of $f$''.
\method Differentiate the claimed answer and simplify to the integrand. Then check the constant separately against an initial condition or a single evaluation point, and check the interval on which the identity holds, since an antiderivative valid on $(0,\infty)$ may be invalid across a singularity.

\begin{worked}
Claim: $\displaystyle\int\sec x\dd x=\ln\abs{\sec x+\tan x}+C$. Differentiating,
\[
\dvop{x}\ln\abs{\sec x+\tan x}=\frac{\sec x\tan x+\sec^{2}x}{\sec x+\tan x}=\frac{\sec x(\tan x+\sec x)}{\sec x+\tan x}=\sec x .
\]
As a numerical cross-check, the central difference of $\ln\abs{\sec x+\tan x}$ at $x=0.7$ is $1.3074593$, and $\sec 0.7=1.3074593$.
\end{worked}

\begin{trap}
Verifying the derivative and forgetting that the constant of integration is constant only on each connected piece of the domain. Across a pole of $\tan x$ the two branches may carry different constants, and an ODE solution assembled from $\ln\abs{y}$ inherits the same defect.
\end{trap}

\begin{seealsob}The independent second method; degenerate and limiting cases.\end{seealsob}

The next check applies when there is nothing to differentiate and no parameter to push: bound the object crudely, and see whether the claimed number is even in the right region.

\tech{Sign and order of magnitude}{core}
\cue Any numerical answer produced by a long chain of substitutions, especially one whose sign changed at some point in the chain.
\method Bound the object with the crudest available inequality. A positive integrand forces a positive integral. A monotone integrand on $[a,b]$ forces a value between $(b-a)f(a)$ and $(b-a)f(b)$. A convex one is bracketed below by the midpoint rectangle and above by the trapezoid. For a sum, compare with the integral of the same expression.

\begin{worked}
Any claimed value for $\displaystyle\int_{0}^{1}\frac{\dd x}{\sqrt{1-x^{4}}}$ must lie in $[1,2]$. Below: the integrand exceeds $1$ on $(0,1)$, so the value exceeds $1$. Above: $1-x^{4}=(1-x)(1+x+x^{2}+x^{3})\geq 1-x$, so the integrand is at most $(1-x)^{-1/2}$, whose integral over $[0,1]$ is $2$. The true value is $1.3110288$, and the two inequalities took ten seconds.
\end{worked}

\begin{trap}
Applying a positivity argument to an integrand that changes sign on the interval. Then the crude bound controls $\int\abs{f}$ only, and the integral itself may be anywhere between $-\int\abs{f}$ and $\int\abs{f}$.
\end{trap}

\begin{seealsob}Numerical evaluation to several digits; squeeze arguments.\end{seealsob}

Symmetry is the check that costs nothing when it applies, because the symmetry was usually already noticed while solving. Its power is that it constrains the whole answer at once rather than one number.

\tech{Symmetry of the answer}{core}
\cue A problem invariant under exchanging two parameters, two variables, or two roles; a definite integral fixed by an involution; a count invariant under relabelling.
\method If the statement is unchanged by a transformation, the answer must be unchanged too. A closed form that is not symmetric under a swap the problem is symmetric under is wrong, without any further computation. The same argument run forwards is often the fastest derivation.

\begin{worked}
Claim: $\displaystyle\int_{0}^{\infty}\frac{\dd x}{(1+x^{2})(1+x^{a})}=\frac{\pi}{4}$ for every real $a$. The substitution $x\mapsto 1/x$ fixes $(0,\infty)$, fixes $\frac{\dd x}{1+x^{2}}$, and sends $a\mapsto -a$; so the value is even in $a$. Adding the two copies gives $\int_{0}^{\infty}\frac{\dd x}{1+x^{2}}=\frac{\pi}{2}$, hence the value is $\frac{\pi}{4}$ independently of $a$. Numerically the integral is $0.7853982$ at $a=3$, at $a=0.7$, and at $a=\pi$. Any claimed answer depending on $a$ dies at the first line.
\end{worked}

\begin{trap}
Asserting a symmetry the problem does not have. Frullani's integral is antisymmetric, not symmetric, under exchanging its two parameters, so its answer $\ln(b/a)$ must change sign under the swap; a proposed answer symmetric in $a$ and $b$ is wrong for the opposite reason.
\end{trap}

\begin{seealsob}Look for the involution; dimensional and scaling sanity.\end{seealsob}

For discrete answers there is a check with no analogue in the continuous world, and it is close to decisive: compute the thing directly for small inputs.

\tech{Small-case brute force}{core}
\cue A count, a probability, the closed form of a recurrence, a combinatorial identity, or any claim beginning ``in general''.
\method Instantiate at $n=0,1,2,3,4$ and compare with direct enumeration. Also check the type constraints: a count must be a non-negative integer, a probability must lie in $[0,1]$, a family of probabilities must sum to $1$, and the growth rate must match intuition. Choose at least one case that was not used in deriving the formula, and at least one degenerate case.

\begin{worked}
Claim: the number of derangements of $n$ objects is $D_{n}=n!\sum_{k=0}^{n}\frac{(-1)^{k}}{k!}$. The formula gives $0,1,2,9,44,265$ for $n=1,\dots,6$; exhaustive enumeration of all $720$ permutations of six objects gives $265$ fixed-point-free ones, and the smaller cases match by hand. A formula that survives $n\leq 4$ against brute force is usually right; one that fails at $n=1$ is always wrong.
\end{worked}

\begin{trap}
Testing only $n=1$ and $n=2$, where almost every plausible wrong formula agrees with the right one, and testing only the case the formula was fitted to.
\end{trap}

\begin{seealsob}Degenerate and limiting cases; hand enumeration with an exhaustiveness statement.\end{seealsob}

The last rung is the most expensive and the only one that is close to conclusive. Two derivations that share no decisive step are two independent measurements of the same quantity.

\tech{The independent second method}{core}
\cue Any answer that matters; any derivation containing a step you were unsure of, especially an interchange of limits, a contour deformation, or a rearrangement.
\method Recompute by a route that shares no step with the first. Reflection versus parameter differentiation; series expansion versus a substitution; a bijection versus a recursion. The point is independence: if both routes rely on the same interchange of summation and integration, they fail together and prove nothing.

\begin{worked}
$\displaystyle\int_{0}^{1}\frac{x-1}{\ln x}\dd x$. Method A, parameter differentiation: with $I(s)=\int_{0}^{1}\frac{x^{s}-1}{\ln x}\dd x$ we get $I'(s)=\int_{0}^{1}x^{s}\dd x=\frac{1}{s+1}$ and $I(0)=0$, so $I(s)=\ln(s+1)$ and $I(1)=\ln 2$. Method B, integral representation: $\frac{x-1}{\ln x}=\int_{0}^{1}x^{t}\dd t$, so exchanging the order gives $\int_{0}^{1}\frac{\dd t}{t+1}=\ln 2$. Numerically the integral is $0.6931472$.
\end{worked}

\begin{trap}
A ``second method'' that reuses the first method's decisive step. Differentiating under the integral sign twice, in two different-looking parametrisations, is one method run twice.
\end{trap}

\begin{seealsob}Numerical evaluation to several digits; symmetry of the answer.\end{seealsob}

\section{What the checks actually catch}

Verification is worth more when you know which errors you are hunting. The six failure modes below account for most wrong answers in calculus, and each has a check on the ladder that is built for it.

\textbf{Bounds not transformed under a substitution.} In $\int_{0}^{1}x\sqrt{1-x^{2}}\dd x$ with $u=1-x^{2}$ and $\dd u=-2x\dd x$, the honest form is $-\frac12\int_{1}^{0}\sqrt{u}\dd u=\frac12\int_{0}^{1}\sqrt{u}\dd u=\frac13$. Keeping the minus sign and the limits $0$ to $1$ together gives $-\frac13$. The sign check kills this in one second, since the integrand is non-negative on $[0,1]$; the same slip on an integrand of mixed sign needs numerical evaluation instead.

\textbf{A lost absolute value in $\ln\abs{u}$.} The antiderivative of $1/x$ is $\ln\abs{x}$, so $\int_{-2}^{-1}\frac{\dd x}{x}=\ln 1-\ln 2=-\ln 2\approx -0.693147$; written as $\ln x$ it is not defined anywhere on the interval. The same applies to $\int\tan x\dd x=-\ln\abs{\cos x}+C$ and to every separable ODE that produces $\ln\abs{y}$. Differentiating the answer back catches it, provided you also ask on which interval the identity holds.

\textbf{A lost solution from dividing by something that can vanish.} Separating $\dv{y}{x}=y^{2}-y$ requires dividing by $y^{2}-y$, which silently discards the constant solutions $y\equiv 0$ and $y\equiv 1$. The general family that emerges is $y=\frac{1}{1-C\eu^{x}}$, which recovers $y\equiv 1$ at $C=0$ but never produces $y\equiv 0$: that solution is genuinely lost. The same slip in trigonometry turns $\sin 2x=\sin x$ into $\cos x=\frac12$, discarding every $x=k\pi$. The check is to substitute each discarded root back into the original equation before dividing.

\textbf{A wrongly chosen branch of an inverse function.} Applying $t=\tan\frac{x}{2}$ to $\int_{0}^{2\pi}\frac{\dd x}{2+\cos x}$ and reading the limits as $t=0$ to $t=0$ gives the answer $0$; the substitution is illegitimate because $t$ has a pole at $x=\pi$ and is not a continuous bijection on the interval. The true value is $\frac{2\pi}{\sqrt3}\approx 3.627599$. The order-of-magnitude check is decisive: the integrand is at least $\frac13$ everywhere, so the integral is at least $\frac{2\pi}{3}\approx 2.094$. The general rule is that $\arctan(\tan\theta)=\theta$ only on $(-\frac{\pi}{2},\frac{\pi}{2})$, and every Weierstrass substitution must be split at the poles of the tangent.

\textbf{A rearranged conditionally convergent series.} The alternating harmonic series sums to $\ln 2\approx 0.693147$, but taking two positive terms for each negative one rearranges it to $\frac32\ln 2\approx 1.039721$, and by Riemann's theorem some rearrangement converges to any prescribed value. The check is to test absolute convergence before rearranging, reordering, or interchanging a sum with an integral, and then to compare partial sums of the series as actually written against the claimed value.

\textbf{A limit computed along one path only.} For $f(x,y)=\frac{xy}{x^{2}+y^{2}}$ the value along either axis is $0$ and the value along $y=x$ is $\frac12$, so the joint limit at the origin does not exist even though both iterated limits are $0$. Two paths is the minimum; the polar form $f=\frac12\sin 2\theta$, independent of $r$, settles it at once. The analogous one-sided error in a single variable is computing a limit as $x\to a^{+}$ and reporting it as the limit.

\section{A wrong answer, caught three ways}

Consider the claim
\[
\int_{0}^{\pi/2}\ln(\sin x)\dd x \;\overset{?}{=}\; -\pi\ln 2 .
\]
This is not a careless answer. It comes from a real and largely correct derivation: the Fourier expansion $\ln\bigl(2\sin\tfrac{\theta}{2}\bigr)=-\sum_{k\geq 1}\frac{\cos k\theta}{k}$ integrates termwise over $[0,\pi]$ to give $\int_{0}^{\pi}\ln\bigl(2\sin\tfrac{\theta}{2}\bigr)\dd\theta=-\sum_{k\geq 1}\frac{\sin k\pi}{k^{2}}=0$, hence $\int_{0}^{\pi}\ln\sin\tfrac{\theta}{2}\dd\theta=-\pi\ln 2$. The solver then substitutes $x=\theta/2$ and forgets that $\dd\theta=2\dd x$. Everything except one factor of two is right, and the answer has the correct sign, the correct constant $\ln 2$, and the correct multiple of $\pi$. Three checks kill it.

\emph{Check one: a crude bound, ten seconds.} On $(0,\tfrac{\pi}{2})$ concavity of the sine puts the graph above its chord, so $\sin x\geq \frac{2x}{\pi}$ and therefore $\ln\sin x\geq\ln\frac{2x}{\pi}$. Integrating the lower bound with $u=2x/\pi$,
\[
\int_{0}^{\pi/2}\ln\sin x\dd x\;\geq\;\frac{\pi}{2}\int_{0}^{1}\ln u\dd u=-\frac{\pi}{2}\approx -1.570796 .
\]
The claimed value $-\pi\ln 2\approx -2.177586$ lies below that bound, so it is impossible. No computation, no numerics, one inequality.

\emph{Check two: numerical evaluation.} The integrand has an integrable logarithmic singularity at $0$, so use the midpoint rule, which never evaluates the endpoint. With a thousand subintervals it returns $-1.08825$; a tanh-sinh quadrature returns $-1.0887930$. The claim is off by a factor of exactly $2$, which is the signature of a lost Jacobian and tells you where to look.

\emph{Check three: an independent second method.} Reflection gives $\int_{0}^{\pi/2}\ln\cos x\dd x=I$, where $I$ is the integral in question. Then
\[
2I=\int_{0}^{\pi/2}\ln(\sin x\cos x)\dd x=\int_{0}^{\pi/2}\ln\frac{\sin 2x}{2}\dd x=-\frac{\pi}{2}\ln 2+\frac12\int_{0}^{\pi}\ln\sin u\dd u,
\]
where the substitution $u=2x$ has been done carefully. Since $\sin u$ is symmetric about $u=\pi/2$, the last integral is $2I$, so $2I=-\frac{\pi}{2}\ln 2+I$ and
\[
I=-\frac{\pi}{2}\ln 2\approx -1.0887930 .
\]
This route uses no Fourier series and no termwise integration, so it fails independently of the first derivation. It agrees with the numerics to seven digits.

\begin{keynote}[A normalisation worth memorising]
The correct value says that $\int_{0}^{\pi/2}\ln(2\sin x)\dd x=0$, and likewise $\int_{0}^{\pi}\ln(2\sin x)\dd x=0$. Any answer in the log-sine family can be tested against this in one line, and the factor of $2$ inside the logarithm is precisely what a lost Jacobian destroys.
\end{keynote}

\section{Recognition matrix}
\begin{recog}{@{}p{0.30\textwidth}p{0.36\textwidth}p{0.26\textwidth}@{}}
\rowcolor{mist}\mh{Answer shape} & \mh{Check that would catch an error} & \mh{What it catches} \\
\midrule
\endhead
A definite integral of a positive integrand & Sign, plus a crude upper bound & Sign slips, lost factors \\
A closed form containing a parameter $a$ & Set $a\to 0$, $a=1$, $a\to\infty$ & Wrong power, missing constant \\
An answer over a scale-invariant domain & Rescale $x=au$ and count powers & Wrong homogeneity degree \\
An indefinite integral & Differentiate it and simplify & Chain-rule factor, lost absolute value \\
An ODE solution & Substitute into the equation and the initial condition & Lost constant solution, wrong $C$ \\
A closed form containing $\ln\abs{u}$ & Test a point where $u<0$ & The dropped absolute value \\
An inverse-trig or Weierstrass result & Evaluate at an interior point and at both endpoints & Wrong branch, pole inside the interval \\
An answer expected symmetric in $m$ and $n$ & Swap $m$ and $n$ & Asymmetric algebra slip \\
An answer claimed independent of $a$ & Evaluate numerically at two values of $a$ & A parameter that did not cancel \\
A count, a probability, or a recurrence & Brute force at $n\leq 5$ & Off-by-one, wrong family \\
A value involving $\pi$, $\ln 2$, $\zeta(2)$ & Numerical agreement to six digits & Rational factors, lost Jacobians \\
A sum of a conditionally convergent series & Partial sums in the original order & Illegitimate rearrangement \\
A limit in more than one variable & A second path, then polar coordinates & Path dependence \\
A limit obtained by expansion & Recompute with one more Taylor order & Premature truncation \\
Anything you are about to submit & Recompute by a technique sharing no step & Everything the others missed \\
\end{recog}
